% Preprint Addendum: Sessions S105-S112
% To be incorporated into preprint.tex (v2.2.0)
% Date: April 2026

\section*{Addendum: Results from Sessions S105--S112}

This addendum records results established after the main preprint was
written. They are to be incorporated in the next arXiv version.

% ============================================================
\subsection*{A1. N=12 Exhaustive Search (Session S35/S112)}
% ============================================================

We extended the N=11 exhaustive search (reported in Section~6) to $N=12$.

\begin{theorem}[N=12 Barrier Confirmation]
No real-valued EML tree with terminals $\{1, x\}$ and $\leq 12$ internal
nodes achieves $\mathrm{MSE} < 0.11$ approximating $\sin(x)$ over
the evaluation domain $[0.1, 1.7]$.
\end{theorem}

\textbf{Search statistics:} 1,704,034,304 trees evaluated; 191,634 passed
the parity filter; best MSE $= 0.1120$; runtime 166 s (8-thread Rust, rayon).
Source: \texttt{monogate-n12/}, results: \texttt{python/results/sin\_n12.json}.

Combined with the analytic proof (T01), this closes the question empirically
to $N=12$ and demonstrates that the barrier is structural, not a matter
of search budget.

% ============================================================
\subsection*{A2. Phantom Attractor: 300-Digit PSLQ Null (Session S105)}
% ============================================================

Using \texttt{mpmath} at 320-digit working precision, we performed PSLQ
on both known phantom attractor basins against a 25-constant basis
$\{1, \pi, e, \ln 2, \ln 3, \ln 5, \sqrt{2}, \sqrt{3}, \sqrt{5}, e^2,
\pi^2, e\pi, \pi/e, \ln\pi, e^\pi, \pi^e, \pi^{1/3}, e^{1/3},
\gamma, \zeta(3), \ln 2 \cdot \ln 3, (\ln 2)^2, G, e^e, \ln\ln 2\}$.

\textbf{Result:} PSLQ NULL at tolerance $10^{-50}$, maxcoeff $= 100$,
for \emph{both} attractor basins:
\begin{itemize}
  \item Primary basin $\alpha_1 \approx 6.21444185277776\ldots$
  \item Secondary basin $\alpha_2 \approx 6.26751862654763\ldots$
\end{itemize}

No minimal polynomial of degree $\leq 8$ with $|\text{coeff}| \leq 500$
was found. \texttt{mpmath.identify()} returns no closed form.

\begin{conjecture}[EML-Generated Transcendence]
The phantom attractor constants $\alpha_1, \alpha_2$ are transcendental
over $\mathbb{Q}$ and lie outside the $\mathbb{Z}$-linear span of the
25 classical constants listed above. They constitute a new class of
\emph{EML-algebraic} constants: values naturally generated by EML
gradient dynamics but not expressible in terms of known transcendentals.
\end{conjecture}

% ============================================================
\subsection*{A3. EML Weierstrass Theorem (Session S109)}
% ============================================================

\begin{theorem}[EML Weierstrass Theorem]
For any compact $K = [a,b]$ with $0 < a < b$, the linear span of all
finite EML trees with leaves $\{1, x\}$ is dense in $C(K)$.
\end{theorem}
\begin{proof}[Proof sketch]
(1) $e^x = \mathrm{eml}(x,1)$ and $\ln x = e - \mathrm{eml}(1,x)$ are exact depth-1 trees.
(2) $x^n = \exp(n\ln x)$ is an exact EML tree for every $n \geq 0$
    (Odrzyw\l{}ek 2026 BEST routing).
(3) All polynomials lie in the EML span (linear combinations of monomials).
(4) Classical Weierstrass: polynomials dense in $C[a,b]$. \qed
\end{proof}

Full proof with lemmas: \texttt{python/paper/eml\_weierstrass\_theorem.tex}.
Lean 4 formalisation in progress: \texttt{lean/EML/InfiniteZerosBarrier.lean}.

\begin{corollary}[Bivariate density]
With leaves $\{1,x,y\}$, the bivariate EML span is dense in $C([a,b]^2)$.
\end{corollary}

% ============================================================
\subsection*{A4. EML Complexity Classification (Session S111)}
% ============================================================

We define a computable complexity measure: $f \in \mathrm{EML}\text{-}k$
if the SVD projection floor $F_f(k) < 10^{-8}$ and $F_f(k-1) \geq 10^{-8}$.

Key results from a 23-function census:
\begin{itemize}
  \item \textbf{EML-1:} $\{e^x, \ln x, 1\}$ — exact depth-1 atoms.
  \item \textbf{EML-2:} $\{x, 1/x, \sqrt{x}, \sinh, \cosh, x/(1+x)\}$.
  \item \textbf{EML-3:} All standard transcendentals: sin, cos, tanh, erf,
        Gaussian $e^{-x^2}$, $x e^{-x}$, $1/(1+x^2)$, $\ln(1+x^2)$, and
        \emph{all} monomials $x^2, x^3, x^4, x^5$.
  \item \textbf{EML-$\infty$:} $\sin x$ on $\mathbb{R}$; $|x|$
        (Infinite Zeros Barrier, non-analyticity).
\end{itemize}

The $N=2 \to N=3$ floor jump ($\sim 10^6\times$ for sin, cos, erf)
is the \emph{N=3 singularity}: the first appearance of algebraically
independent (EL-height-3) directions in the atom dictionary.

Full census: \texttt{python/paper/eml\_complexity\_census.tex}.

% ============================================================
\subsection*{A5. Grammar G3 and Negative-Exponent Barrier (Session S110)}
% ============================================================

Grammar G3 (EML + DEML + EXL) was tested on 9 target functions.
All 9 are G3-DENSE (MSE $< 10^{-8}$) including $e^{-x}$, $e^{-x^2}$,
and $x e^{-x}$ which require the DEML gate.
The G1 (EML-only) span is already DENSE for all these functions at depth 3.

\textbf{Conclusion:} At depth $\geq 3$ the DEML gate provides no barrier-lifting
advantage for smooth targets — G1 and G3 are equally dense.
The DEML gate's value is in \emph{node efficiency}: $e^{-x}$ as a 1-node
DEML primitive vs. the EML-alone approximation which needs depth $\geq 3$.
For piecewise-smooth targets ($|sin|$, triangle wave, sign function),
G2 (EML+DEML) gives 6--16$\times$ improvement but sawtooth/sign remain open.

% ============================================================
\subsection*{A6. Complex EML Depth-7 Analysis (Session S109)}
% ============================================================

The depth-6 near-miss $\mathrm{Im}(\text{EML}_1) \approx 0.9999952$
(gap $4.76 \times 10^{-6}$, found in session S93) was further analysed.

Two routes to $\mathrm{Im} = 1$ from the grammar $\{1\}$:
\begin{enumerate}
  \item Via $\pi/\tan(1) \approx 2.01719$: gap $4.19 \times 10^{-5}$ at depth 6.
  \item Via $\pi/2 \approx 1.5708$: gap $9.75 \times 10^{-6}$ at depth 6.
\end{enumerate}

Both gaps are shrinking with depth (density conjecture), but neither is
zero.  By Hermite--Lindemann, $\tan(1)$ is transcendental, and $\pi/2$
is transcendental.  Whether either lies in the real EML closure from $\{1\}$
is equivalent to the open problem: \emph{is $i$ constructible from $\{1\}$
via EML?}

% ============================================================
\subsection*{A7. Quantum EML Weierstrass (Session S111)}
% ============================================================

44 matrix-EML (meml) depth-1 atoms were evaluated on the space of
$2 \times 2$ density matrices. The atom matrix has rank 4 in the
8-dimensional real matrix space. All 6 tested density matrices
(pure $X,Y,Z$ eigenstates, mixed states, Bell states) were approximated
with MSE $< 3 \times 10^{-32}$ (machine precision).

\begin{conjecture}[Quantum EML Weierstrass]
For $d \times d$ density matrices, there exists a finite set of
depth-$k(d)$ meml atoms that spans the full $(d^2-1)$-dimensional
density matrix manifold. For $d=2$, $k=1$ suffices.
\end{conjecture}

% ============================================================
\subsection*{A8. Pumping Lemma: Tight Bound Conjecture (Session S107)}
% ============================================================

The formal Pumping Lemma states: an EML tree of depth $k$ has at most
$2^k$ zeros on any compact interval (proof by repeated application of
Rolle's theorem to the nested exp/log structure).

Empirical survey of representative depth-$k$ trees on $[-8,8]$ shows
the observed maximum is 0--1 zeros per tree, far below $2^k$.

\begin{conjecture}[Tight Pumping Bound]
The tight bound is $O(k)$, not $O(2^k)$.
EML trees at depth $k$ are almost always monotone (0 zeros)
because exp and log both map intervals to intervals without oscillation.
The $2^k$ bound is an artefact of the Rolle argument, not a tight count.
\end{conjecture}

This is consistent with the fact that all EML-$k$ trees (finite depth)
are real-analytic and have isolated zeros, while the \emph{growth} in zero
count is slow.

% ============================================================
\subsection*{Updated Open Problems}
% ============================================================

\begin{enumerate}
  \item \textbf{Is $i$ constructible from $\{1\}$ via EML?}
    Equivalently: is $\pi/2$ or $\pi/\tan(1)$ in the real EML closure?
    (Complex gap $4.76 \times 10^{-6}$ at depth 6; density conjecture would
    imply yes in the limit.)
  \item \textbf{Are the phantom attractors transcendental over the EL field?}
    PSLQ null at 320 digits, no minimal poly of degree $\leq 8$.
    Conjecture: yes, they form a new EML-algebraic constant class.
  \item \textbf{Formal proof of Lemma 2.}
    The algebraic proof that $x^n$ is an \emph{exact} EML tree
    (not just approximated) is implicit in the BEST routing construction.
    A self-contained Lean 4 proof is the target.
  \item \textbf{Tight Pumping Lemma bound.}
    Is the zero-count bound $O(k)$ or $O(2^k)$?
    Constructing a depth-$k$ EML tree with $\Omega(k)$ zeros would resolve this.
  \item \textbf{G2 density for sawtooth/sign.}
    Grammar G2 (EML+DEML) improves $|sin|, |cos|$, triangle by 6--16$\times$
    but sawtooth and sign function remain open.
    G3 adds no benefit over G1 for smooth functions at depth $\geq 3$.
\end{enumerate}

\vspace{1em}
\noindent\textit{Version note: This addendum covers Sessions S105--S112 (April 2026).
The preprint will be updated to v2.2.0 on the next arXiv revision.}
